
\section{Frequently asked questions}
\label{sec:trouble}


\begin{itemize}
    \item {\bf Q:} I get an error message "workflow not found",  and/or nothing happens after I submitted the {\tt edps} command with my workflow. How do I fix it? \newline
    {\bf A:}  This should not happen if the recommended installation procedures for edps and pipelines have been followed. If the installation was done in a non-standard way, the "workflow not found" message could indicate that:
    \begin{itemize}
        \item the {\tt esorex} command is not associated to the instrument pipeline you want to use. 
        \item The installed pipeline does not have an EDPS workflow yet.
        \item The workflow name is mispelled.
     \end{itemize}
  
        Type command \verb|esorex --recipes| to check which pipeline is seen by {\tt esorex}. If there are no recipes for your instrument, it means that path in {\tt application.properties} points to the wrong {\tt esorex} installation. Fix it and restart the edps server. Type the command {\tt edps -lw} to list the installed workflows. If the workflow is not present, please check the spelling or the workflow directory in the {\tt application.properties} file. See Section \ref{sec:setup}


    \item {\bf Q:} The association reveals that my datasets are not complete, but I think to have all the needed data. How do I fix it?\newline
    {\bf A:} It could be that EDPS could not find the location of static calibrations for that instrument pipeline (this could happen, for example, if the instrument pipeline was not installed following the recommended procedures). Try to add the static calibration directory to the list of input data via the {\tt -i} option.
      \item {\bf Q:} I do not remember the name of the workflow I want to run. How do I know the exact workflow names? 
      \newline
    {\bf A:} EDPS workflows are installed together with the instrument pipeline installation. To see which are the workflows installed in your system and their names, type: {\tt edps -lw}. 
    \item {\bf Q:} I edited the configuration file application.properties, but this seems to have no effect. Why?\\
    {\bf A:} In order for the changes to take effect, first close the EDPS server by typing the command {\tt edps -shutdown}, and then relaunch the request.
    \item{\bf Q:} How do I know which recipe parameters were or will be used?
    \newline{\bf A:} The book-keeping directory has the file parameters.rc, which indicates only the values of the recipe parameters that are different from the recipe default. The products have record of the recipe parameters values in their header. For inspecting the parameters values that are going to be used in the reduction {\it before} the reduction starts, consult Section \ref{sec:edps_default_params}.
    \item {\bf Q:} The {\tt edps} command cannot connect, or I started the server on a non-default port. What should I do?
    \newline{\bf A:} The client defaults to {\tt localhost:5000}. Pass {\tt -P <port>} (and {\tt -H <host>} if needed). After editing {\tt application.properties}, always {\tt edps --shutdown} before the next request.
\end{itemize}